\section{Programação Estocástica Multiestágio}
\label{sec:prog_estocastica}

A Programação Estocástica (PE) é o ramo da otimização matemática que trata problemas de decisão em que parte dos parâmetros do modelo é incerta no momento da decisão. A incerteza é representada por variáveis aleatórias, e o objetivo é encontrar uma política que seja ótima em valor esperado — ou em outro critério de risco — sobre a distribuição dessas variáveis \cite{birge2011}.

\subsection{Formulação Multiestágio e Não Antecipatividade}
\label{subsec:multistagio}

Em muitos problemas práticos, as decisões são tomadas sequencialmente e as incertezas se revelam de forma progressiva. A Programação Estocástica Multiestágio (PEM) modela essa estrutura temporal: em cada estágio $t$, a decisão é tomada com base nas informações disponíveis até aquele momento, sem acesso a realizações futuras das variáveis aleatórias. Essa propriedade é denominada \emph{não antecipatividade} e é uma restrição fundamental da PEM \cite{shapiro2009}.

A incerteza é tipicamente representada por uma árvore de cenários, na qual cada nó corresponde a uma realização observável das variáveis aleatórias e cada ramo representa uma transição probabilística entre estágios. A formulação determinística equivalente (DEQ) constrói explicitamente essa árvore e resolve o problema resultante como um único programa linear. O tamanho do DEQ cresce exponencialmente com o número de estágios e ramificações, tornando-o inviável para instâncias com horizontes longos ou alta discretização das incertezas \cite{birge2011}.

\subsection{Técnicas de Decomposição}
\label{subsec:decomposicao}

Para contornar a explosão combinatória da DEQ, técnicas de decomposição exploram a estrutura temporal do problema para resolvê-lo por estágios. A decomposição de Benders e suas variantes são as abordagens mais utilizadas nesse contexto \cite{shapiro2009}. A Programação Dinâmica Dual Estocástica, descrita na Seção~\ref{sec:pdde_literatura}, é a técnica de decomposição adotada neste trabalho.
